\exnumber{03}
\extitle{UX V - Ex 4 - Tri de cartes}
\exfiles{ex03/archive zip, nomme Name\_FirstName\_DeliverableName\_Date\_VersionNumber}
\exauthorize{Ceux nommés dans l'exercice}
\makeheaderfiles

\textbf{Instructions :}

Dans la situation actuelle, les sites web ne sont pas idéaux pour accomplir les Jobs-to-be-Done que votre utilisateur a en tête. Alors pratiquons le card sorting pour réorganiser le contenu du site !

Organisez un atelier de card sorting avec le contenu du site web que vous avez choisi dans l'Ex. 1.

Pour cela :

\begin{itemize}
\item Définir vos objectifs de recherche
\item Préparer le matériel
\item Trouver des participants
\item Animer un atelier de card sorting auprès de 5 personnes
\item Proposer une nouvelle taxonomie (organisation du contenu)
\end{itemize}

\textbf{Liste de sites :}

\begin{itemize}
\item Universal Music Latin America \url{https://www.universalmusica.com/}
\item Delta Records \url{https://deltarecords.net/}
\item Zoho Music \url{https://www.zohomusic.com/}
\item Luaka Bop \url{https://www.luakabop.com/}
\item Rimas Music \url{https://rimasmusic.com/}
\item D'autres... tant que c'est un label de musique sud-américain ;-)
\end{itemize}

\textbf{Livrables :}

\begin{itemize}
\item Un script d'atelier de tri de cartes avec les objectifs de recherche (document PDF)
\item Des photos de l'atelier en cours
\item Un document pour montrer la nouvelle organisation du contenu (ce peut être une photo du board avec des post-its ou la version numérique).
\item Déposez les documents sur Github, dans un zip nommé « ProjectTitle\_Name\_FirstName », et chaque livrable ainsi : « Name\_FirstName\_DeliverableName\_Date\_VersionNumber ». Exemple : « Doe\_John\_MidFiPrototype\_05242024\_V1 ».
\item Ajoutez un titre dans tous les documents écrits.
\item Soignez l'aspect global des livrables. Ils doivent être clairs, simples et lisibles. Vous pouvez regarder canva.com, sans surcharger le dessin de détails.
\end{itemize}

\textbf{Ressources :}

\begin{itemize}
\item Bases pour mener un card sorting \url{https://uxdesign.cc/basics-on-how-to-conduct-card-sorting-6e4ccc7cc261?sk=cafd49e16c71a72c3fad93e63c1998f1}
\item Card Sorting  ---  quoi, comment et les avantages \url{https://uxdesign.cc/card-sorting-what-how-the-perks-29f6cb020270}
\item 10 choses à savoir sur le card sorting \url{http://www.measuringu.com/blog/card-sorting.php}
\item Card sorting : un guide définitif \url{http://boxesandarrows.com/card-sorting-a-definitive-guide/}
\end{itemize}

